\title{Direct Preference Optimization: Your Language Model is Secretly a Reward Model}

\begin{abstract}
Large-scale unsupervised language models (LMs) acquire extensive general knowledge and exhibit some reasoning capabilities, but exerting fine-grained control over their outputs remains a challenge due to the fully unsupervised methodology behind their training. To address this, current approaches typically obtain human evaluations comparing the quality of generated outputs and adapt the LM accordingly, most often through reinforcement learning from human feedback (RLHF). Despite its effectiveness, RLHF presents practical difficulties, including instability and procedural complexity: it involves first constructing a reward model that embodies human preferences, followed by reinforcement learning to guide the LM towards higher reward while preventing excessive deviation from its pretrained state. In this work, we propose a novel way of expressing the reward model in RLHF that admits a closed-form solution for the associated optimal policy, thereby allowing the original RLHF optimization to be reduced to a straightforward classification loss. This method, termed \textit{Direct Preference Optimization} (DPO), is robust, efficient, and computationally streamlined, removing the necessity for language model sampling during fine-tuning and minimizing requirements for hyperparameter searches. Empirical results demonstrate that DPO aligns LMs to human preferences on par with or better than prevalent techniques. In particular, DPO surpasses RLHF approaches based on PPO in modulating generative sentiment and achieves equivalent or superior results in summarization and single-turn dialogue tasks, all while simplifying the fine-tuning process considerably.
\end{abstract}

\section{Introduction}
Large-scale unsupervised language models (LMs), when trained on extensive corpora, exhibit a range of remarkable abilities~\citep{chowdhery2022palm, brown2020language, touvron2023llama, bubeck2023sparks}. Nonetheless, because these models are exposed to text generated by humans with divergent intentions, objectives, and expertise, they also absorb patterns that may not be desirable to emulate. For instance, while we may want an AI assistant for programming to \textit{recognize} typical programming errors to facilitate their correction, our preference would be for the model to generate code that reflects the (sometimes infrequent) instances of high-quality code from its training set. In a similar vein, it is useful for the language model to be \textit{familiar} with a widespread misconception shared by 50\% of individuals, but we certainly do not want it to affirm this misconception as correct half the time it is queried. Thus, the ability to control a model’s \emph{preferred outputs and conduct}—drawing selectively from its broad \textit{knowledge and capabilities}—is essential for creating AI systems that are safe, effective, and manageable \citep{ouyang2022training}. Although present-day techniques generally guide LMs to reflect human preferences through reinforcement learning (RL), we demonstrate that the RL-driven objective currently employed can, in fact, be solved precisely with a straightforward binary cross-entropy objective, thereby streamlining the preference learning process.

In general terms, mainstream approaches shape model behavior by leveraging carefully constructed datasets of human preferences that encode what is considered safe and beneficial. This stage of preference learning follows an initial phase of large-scale unsupervised pre-training on text. Although the most direct method for incorporating preferences is to fine-tune the model on expert demonstrations—supervised learning of high-quality examples—the leading paradigm is reinforcement learning from human (or AI) feedback (RLHF/RLAIF; \citep{christiano2017deep,bai2022constitutional}). RLHF methods involve training a reward model on annotated human preference data, and subsequently using RL to adapt the language model to generate responses which the reward model ranks highly, while preventing significant divergence from its original outputs. Despite yielding language models that excel in conversation and programming tasks, RLHF requires a more elaborate training regime compared to standard supervised approaches, involving the training of several models and repeated policy sampling during the training process, resulting in considerable computational overhead.

This work presents a method for directly aligning language models with human preferences, sidestepping the need for explicit reward modeling or reinforcement learning approaches. We introduce \textit{Direct Preference Optimization} (DPO), an algorithm that tacitly targets the same objective as prevalent RLHF methods—maximizing rewards under a KL-divergence regularization—while offering ease of implementation and training. Conceptually, DPO modifies the relative log-likelihood of preferred versus non-preferred responses and employs adaptive, instance-specific importance weights to counteract the degradation issues observed with straightforward probability ratio-based objectives. Consistent with established frameworks, DPO utilizes a theoretical preference model (e.g., the Bradley-Terry model; \cite{bradley1952rankanalysis}) to evaluate how closely a reward function matches real-world preference judgments. Unlike standard approaches that utilize the preference model to inform a preference-based loss for training a reward model, followed by policy optimization with this model, DPO applies a variable transformation so that the preference loss is a direct function of the policy itself. With access to a corpus of human-annotated preferences over model outputs, DPO can train a policy using a basic binary cross entropy loss, yielding the optimal policy for a latent reward function tailored to the collected preference data.

The principal contribution of this work is Direct Preference Optimization (DPO), an uncomplicated algorithm that obviates reinforcement learning in preference-based training of language models. Experimental evaluations reveal that DPO performs on par with, or better than, existing techniques—including RLHF methods utilizing PPO—on a range of preference-learning tasks such as sentiment steering, summarization, and conversational modeling, using models containing up to 6B parameters.

\section{Related Work}

Recent advancements in self-supervised language models have demonstrated that, as model scale increases, these systems are capable of performing certain tasks in a zero-shot manner \citep{radford2019language} or with the help of few-shot prompting \citep{gpt3,megatron,chowdhery2022palm}. Nevertheless, their efficacy on downstream applications and their conformity to user intentions are notably enhanced through fine-tuning with datasets that pair instructions with human-generated completions \citep{mishra-etal-2022-cross,sanh2022multitask,chung2022scaling,thoppilan2022lamda}. This process, commonly termed `instruction-tuning,' facilitates the capacity of large language models to extrapolate to novel instructions not present in the fine-tuning corpus, thereby broadening their overall applicability \citep{chung2022scaling}. Although instruction tuning has proven effective, collecting human comparative assessments of response quality is often more practical than acquiring expert-curated demonstrations. In light of this, more recent work has pivoted to fine-tuning LLMs using datasets constructed from human preference judgments, thereby advancing performance on tasks such as translation \citep{kreutzer-etal-2018-reliability}, summarization \citep{stiennon2022learning,ziegler2020finetuning}, story generation \citep{ziegler2020finetuning}, and following user instructions \citep{ouyang2022training,ramamurthy2023is}. These approaches initially involve training a neural reward model that aligns with observed human preferences, frequently modeled by frameworks such as the Bradley-Terry model \citep{bradley1952rankanalysis}. The resulting reward signal is then used to further train the language model via reinforcement learning algorithms—including REINFORCE \citep{williams1992reinforce}, proximal policy optimization (PPO; \cite{schulman2017proximal}), or other similar techniques \citep{ramamurthy2023is}. In parallel, related research exploits instruction-following LLMs refined through human feedback to synthetically generate new preference data targeting attributes like safety or harmlessness \citep{bai2022constitutional}, relying on weak supervision based on rubric-style textual feedback from humans for the model's labeling process. This body of work bridges research on reinforcement learning for language modeling across various goals~\citep{Ranzato2015SequenceLT,paulus2018a,wu2018learning} and more general approaches for training with human preferences \citep{christiano2017deep,kupcsik2018learning}. Nonetheless, even with the advantages of relative human preference data, applying reinforcement learning for LLM fine-tuning at scale remains technically challenging; the present work introduces a theoretically grounded method for optimizing based on relative preferences that avoids the complexities of RL.

Research on learning policies from preferences outside the language domain has been conducted within both bandit and reinforcement learning frameworks, leading to the development of several methods. When contextual bandit algorithms leverage preferences or rankings over actions instead of explicit reward signals, this paradigm is termed the contextual dueling bandit (CDB; \cite{yue2012karmed,dudik2015contextual}). In scenarios without access to absolute rewards, CDB theory replaces the optimal policy with the concept of a \textit{von Neumann winner}—that is, a policy whose expected probability of outperforming any alternative policy meets or exceeds 50\% \citep{dudik2015contextual}. Notably, preference information in CDBs is acquired in an online fashion, whereas human preference learning typically relies on a static set of previously collected, preference-labeled action pairs \citep{yan2022human}. In the realm of \textit{preference-based RL} (PbRL), binary preferences are sourced from an opaque `scoring' function rather than direct reward feedback \citep{BusaFekete2014,ruiz2023dueling}. Multiple approaches for PbRL exist—including those that make use of off-policy preference datasets—but most share the characteristic of explicitly learning the underlying scoring (reward) function prior to policy optimization \citep{jain2013learning,BusaFekete2014,christiano2017deep,sadigh2017active,kupcsik2018learning}. Our approach deviates from this tradition by employing a single-stage methodology that directly optimizes a policy with respect to given preferences.

\section{Preliminaries}\label{section:prelims}

Here, we summarize the RLHF pipeline as described by \citeauthor{ziegler2020finetuning} and subsequent works \citep{stiennon2022learning, bai2022training, ouyang2022training}. This process is typically delineated into three stages: 1) supervised fine-tuning (SFT); 2) the joint step of collecting preference data and reward modelling; and 3) reinforcement learning-based optimization.

\textbf{SFT}: The standard RLHF workflow commences by adapting a pre-trained language model through supervised learning on high-quality datasets relevant to the downstream application (e.g., dialogue, summarization), producing a fine-tuned model denoted as $\pi^\text{SFT}$.

\textbf{Reward Modelling Phase}: The next phase utilizes the SFT model, which is given input prompts $x$ and generates answer pairs $(y_1, y_2)\sim \pi^\text{SFT}(y \mid x)$. Human annotators are then asked to express a preference between the two responses; this is denoted $y_w\succ y_l \mid x$, with $y_w$ as the chosen and $y_l$ as the rejected response from the pair $(y_1, y_2)$. These preferences are presumed to originate from a hidden reward function $r^*(y, x)$, the details of which remain unknown. A range of models exist for representing such preferences, among which the Bradley-Terry (BT) model \cite{bradley1952rankanalysis} is widely adopted, though the more general Plackett-Luce model \citep{plackett1975analysis, luce2012individual} can also be employed if multi-way rankings are provided. The BT model expresses the human preference distribution $p^*$ as follows:

\begin{equation}\label{eq:bradley-terry}
    p^*(y_1\succ y_2 \mid x)=\frac{\exp\left(r^*(x, y_1)\right)}{\exp\left(r^*(x, y_1)\right) + \exp\left(r^*(x, y_2)\right)}.
\end{equation}

Given a static dataset of human comparisons $\mathcal{D} = \bigl\{x^{(i)}, y_w^{(i)}, y_l^{(i)}\bigr\}_{i=1}^N$ drawn according to $p^*$, we introduce a parametric reward function $r_{\phi}(x, y)$ and learn its parameters by maximizing likelihood. This preference learning setup can be reframed as a binary classification task, leading to the following negative log-likelihood loss:

\begin{equation}\label{eq:reward_model}
    \mathcal{L}_R(r_{\phi}, \mathcal{D}) = -\mathbb{E}_{(x, y_w, y_l)\sim \mathcal{D}}\bigl[\log \sigma(r_{\phi}(x, y_w)- r_{\phi}(x, y_l))\bigr]
\end{equation}

where $\sigma$ denotes the logistic sigmoid. For language models, $r_{\phi}(x, y)$ is typically initialized from a supervised fine-tuned (SFT) model $\pi^\text{SFT}(y \mid x)$, with an additional linear output layer added atop the transformer’s final representation to produce a single scalar reward estimate \cite{ziegler2020finetuning}. To mitigate reward estimate variance, it is common practice to normalize rewards, enforcing $\mathbb{E}_{x,y\sim \mathcal{D}}\left[r_\phi(x, y)\right] = 0$ for all $x$.

\textbf{RL Fine-Tuning Phase}: In the subsequent RL fine-tuning stage, this learned reward provides feedback for further optimizing the language model. Following prior approaches~\citep{jaques2017sequence, jaques2020human}, the objective is formalized as:

\begin{equation}\label{eq:RL}
\max_{\pi_{\theta}}  \mathbb{E}_{x\sim \mathcal{D}, y\sim \pi_{\theta}(y \mid x)}\bigl[r_{\phi}(x, y)\bigr] - \beta\mathbb{D}_{\textrm{KL}}\bigl[\pi_{\theta}(y\mid x)\mid \mid \pi_\text{ref}(y\mid x)\bigr],
\end{equation}

where the regularization parameter $\beta$ controls divergence from the reference policy $\pi_\text{ref}$, usually instantiated as the original SFT model $\pi^\text{SFT}$. In practical settings, $\pi_\theta$ is initialized to $\pi^\text{SFT}$ as well. The KL constraint is critical for ensuring the learned policy stays close to regions where the reward model is valid, preserves output diversity, and avoids collapse onto a few highly rewarded responses. Since language output is discrete, the gradient is unavailable, and optimization generally proceeds with reinforcement learning. Conventional strategies \citep{ziegler2020finetuning, stiennon2022learning, bai2022training, ouyang2022training} define the reward as ${r(x, y) = r_{\phi}(x, y) -\beta (\log \pi_{\theta}(y\mid x) - \log \pi_\text{ref}(y\mid x))}$ and maximize this using PPO \cite{schulman2017proximal}.

\section{Direct Preference Optimization}\label{sec:DPO}

Motivated by the difficulties encountered when applying large-scale reinforcement learning to language model tuning, we aim to develop a streamlined approach for direct policy optimization grounded in preferences. In contrast to conventional RLHF methodologies that first fit a reward model and subsequently optimize it with RL, our method utilizes a specific form of reward model parameterization, allowing us to analytically recover the corresponding optimal policy, thereby sidestepping the need for reinforcement learning procedures. As will be elaborated below, the core insight lies in employing an explicit mapping from the reward function to the optimal policy, permitting a reformulation of the learning objective from the space of reward models to the space of policies. This transformation eliminates the necessity of learning an explicit, standalone reward predictor, while retaining the ability to optimize according to structured human feedback models like the Bradley-Terry formulation. Essentially, the policy network is tasked with representing both the generative language model and the implicit underlying reward function

\textbf{Derivation of the DPO Objective.} We begin with the standard reinforcement learning objective given in Eq.~\ref{eq:RL}, which is defined over a general reward function $r$. Building on established results in the literature~\citep{peters2007reinforcement, peng2019advantage, korbak2022reinforcement, go2023aligning}, one can readily verify that the optimal policy for the reward maximization problem subject to a KL divergence constraint in Eq.~\ref{eq:RL} assumes the following form:

\begin{equation}\label{eq:op_policy}
    \pi_r(y\mid x) = \frac{1}{Z(x)}\pi_\text{ref}(y\mid x)\exp\left(\frac{1}{\beta}r(x, y)\right),
\end{equation}

where $Z(x) =\sum_{y}\pi_\text{ref}(y\mid x)\exp\left(\frac{1}{\beta}r(x, y)\right)$ represents the partition function. The full derivation is provided in Appendix \ref{app:derivation1}. Even if the reward function $r^*$ is estimated using a maximum likelihood approach $r_{\phi}$, the computation of the normalization constant $Z(x)$ remains a computational bottleneck~\citep{korbak2022reinforcement, go2023aligning}, which poses significant challenges in practical scenarios. To circumvent this, we manipulate Eq.~\ref{eq:op_policy} to express the reward function as a function of the corresponding optimal policy $\pi_r$, the reference policy $\pi_\text{ref}$, and the normalization $Z(x)$. By taking logarithms on both sides of Eq.~\ref{eq:op_policy} and rearranging, we arrive at:

\begin{equation}\label{eq:main_eq}
    r(x,y) =\beta \log \frac{\pi_r(y\mid x)}{\pi_\text{ref}(y\mid x)} + \beta \log Z(x).
\end{equation}

This reparameterization is equally applicable to the true reward $r^*$ and its associated optimal policy $\pi^*$. Notably, in the Bradley-Terry framework, only the difference in rewards between two candidates is relevant; that is, ${p^*(y_1 \succ y_2 \mid x) = \sigma(r^*(x, y_1) - r^*(x, y_2))}$. Substituting the reparameterized reward expression from Eq.~\ref{eq:main_eq} into the Bradley-Terry preference equation (Eq.~\ref{eq:bradley-terry}) leads to the cancellation of the partition functions, yielding an expression for human preference probability that depends solely on $\pi^*$ and $\pi_\text{ref}$. Accordingly, under the Bradley-Terry preference model, the RLHF-optimal policy $\pi^*$ satisfies:

\begin{equation}\label{eq:objective}
    p^*(y_1\succ y_2 \mid x)=\frac{1}{1 + \exp\left(\beta \log \frac{\pi^*(y_2\mid x)}{\pi_\text{ref}(y_2\mid x)} - \beta \log \frac{\pi^*(y_1\mid x)}{\pi_\text{ref}(y_1\mid x)}\right)}
\end{equation}

Details regarding this derivation can be found in Appendix~\ref{app:derivation2}. While the above utilizes the Bradley-Terry model, analogous results for more general Plackett-Luce models~\citep{plackett1975analysis, luce2012individual} are derived in Appendix~\ref{app:plackett_luce_models}.

Having expressed the human preference probability purely in terms of the optimal policy, we are able to define a maximum likelihood estimation framework for a parameterized policy $\pi_\theta$. In direct analogy with reward modeling approaches (cf. Eq.~\ref{eq:reward_model}), we propose the following policy learning objective:

\begin{equation}\label{eq:optimum_model}
    \mathcal{L}_\text{DPO}(\pi_{\theta}; \pi_\text{ref}) = -\mathbb{E}_{(x, y_w, y_l)\sim \mathcal{D}}\left[\log \sigma \left(\beta \log \frac{\pi_{\theta}(y_w\mid x)}{\pi_\text{ref}(y_w\mid x)} - \beta \log \frac{\pi_{\theta}(y_l\mid x)}{\pi_\

By adopting this approach, we model an implicit reward through an alternative parameterization, where the corresponding optimal policy coincides with $\pi_\theta$. Notably, the method aligns with optimizing a reparametrized Bradley-Terry model, and consequently inherits beneficial theoretical guarantees such as consistency, provided the preference data distribution satisfies certain conditions \cite{bong2022generalized}. Additional theoretical insights into DPO, especially in comparison to related methodologies, are explored in Section~\ref{sec:theory}.

\textbf{Mechanism of the DPO Update.} To gain a mechanistic perspective on DPO, we examine the gradient of the loss function $\mathcal{L}_\text{DPO}$. The gradient with respect to $\theta$ is expressed as follows:

\begin{multline*}\label{eq:gradient}
    \nabla_\theta \mathcal{L}_\text{DPO}(\pi_\theta;\pi_\text{ref}) = \\ -\beta\mathbb{E}_{(x, y_w, y_l) \sim \mathcal{D}} \bigg[\underbrace{\sigma(\hat{r}_\theta(x, y_l) - \hat{r}_\theta (x, y_w))}_\text{larger impact when reward prediction errs}\bigg[\underbrace{\nabla_\theta\log \pi(y_w \mid x)}_\text{boost probability of $y_w$} - \underbrace{\nabla_\theta\log\pi(y_l \mid x)}_\text{reduce probability of $y_l$}\bigg]\bigg],
\end{multline*}

where $\hat{r}_\theta(x, y) = \beta \log \frac{\pi_\theta(y \mid x)}{\pi_\text{ref}(y \mid x)}$ defines the implicit reward arising from the language model $\pi_\theta$ in relation to the reference model $\pi_\text{ref}$ (further discussed in Section~\ref{sec:theory}). Intuitively, the loss function’s gradient, $\mathcal{L}_\text{DPO}$, promotes preferred responses $y_w$ by increasing their likelihood and discourages dispreferred responses $y_l$ by decreasing theirs. Crucially, the gradient's influence is modulated by the extent to which the implicit reward model $\hat{r}_\theta$ overvalues dispreferred completions, scaled by $\beta$—effectively weighting the update by the degree of misranking, while taking into account the KL regularization. Empirical findings indicate this weighting mechanism is essential; omitting the coefficient—as in a simplistic version of the approach—can lead the language model to collapse, as demonstrated in Appendix Table~\ref{tab:unlikelihood_generations}.

\textbf{DPO framework.} The standard workflow for DPO comprises the following steps: 1) For each prompt $x$, draw sample completions $y_1, y_2 \sim \pi_\text{ref}(\cdot \mid x)$, annotate these with human preference judgments to build the preference dataset $\mathcal{D} = \{x^{(i)}, y_w^{(i)}, y_l^{(i)}\}_{i=1}^N$; 2) update the language model $\pi_\theta$ by minimizing $\mathcal{L}_\text{DPO}$, conditioned on the reference model $\pi_\text{ref}$, dataset $\mathcal{D}$, and chosen $\beta$ parameter. In realistic scenarios, rather than generating new samples and human annotations, one typically relies on existing publicly available preference datasets. Since these datasets originate from samples generated via $\pi^\text{SFT}$, it is standard practice to set $\pi_\text{ref} = \pi^\text{SFT}$ whenever this model is accessible. If, however, $\pi^\text{SFT}$ is unavailable, $\pi_\text{ref}$ is obtained by maximizing the likelihood over preferred responses ${(x, y_w)}$, formally, ${\pi_\text{ref} = \argmax_{\pi}\mathbb{E}_{x, y_w \sim \mathcal{D}}\left[\log \pi(y_w \mid x)\right]}$. This approach serves to bridge the gap between the intractable true reference distribution and the empirical $\pi_\text{ref}$ deployed in DPO. Implementation specifics and hyperparameter choices are described in Appendix~\ref{app:implementation}.

\section{Theoretical Analysis of DPO} This section delves into the theoretical underpinnings of DPO, offering both conceptual clarification and rigorous support, and draws connections between the strengths of DPO and the limitations inherent to actor-critic-based RLHF algorithms (such as PPO~\cite{schulman2017proximal}).

\label{sec:theory}

\subsection{Your Language Model Is Secretly a Reward Model} DPO eliminates the need to explicitly fit a reward model or utilize RL for policy training by employing a unified maximum likelihood framework. It is important to observe that the DPO optimization (see Eq. \ref{eq:main_eq}) aligns with the Bradley-Terry model when rewards are defined as $r^*(x, y) = \beta \log\frac{\pi^*_\theta(y \mid x)}{\pi_\text{ref}(y \mid x)}$. Here, the model $\pi_{\theta}$ is optimized in the same manner as the reward model described in Eq. \ref{eq:reward_model}, given a variable transformation. This section develops the theoretical justification for this reparameterization, demonstrates that it imposes no restrictions on the expressiveness of reward functions that can be learned, and shows that it enables recovery of the optimal policy without approximation. We begin by formalizing an equivalence relation on reward functions.

\begin{definition}
Two reward functions $r(x, y)$ and $r'(x, y)$ are defined as equivalent if and only if there exists a function $f$ such that ${r(x, y)-r'(x, y) = f(x)}$.
\end{definition}

This relation clearly satisfies the properties of an equivalence relation, thereby dividing the space of reward functions into equivalence classes. This leads to two key results:

\begin{lemma}\label{lemma:same_prefrence} Within the framework of Plackett-Luce—specifically including the Bradley-Terry model—any pair of reward functions belonging to the same equivalence class generate identical preference distributions.
\end{lemma}

\begin{lemma}\label{lemma:same_policy}
    For the constrained RL setting, any two reward functions from the same equivalence class yield the same optimal policy.
\end{lemma}

The arguments supporting these statements are elementary and can be found in Appendix \ref{app:lemma1}. The first lemma highlights the familiar issue of under-determination present in Plackett-Luce-type models \cite{plackett1975analysis}. To mitigate this lack of identifiability, additional constraints are typically enforced to ensure that the maximum likelihood estimates derived from Eq. \ref{eq:reward_model} are well-defined \cite{bong2022generalized}. The second lemma implies that, in terms of optimal policy derivation, it is sufficient to recover any member from the target equivalence class of rewards. The following theorem, proved in Appendix~\ref{app:thm1}, formalizes this characterization:

\begin{theorem}\label{thm:main}
    Subject to mild conditions, every reward equivalence class consistent with Plackett-Luce (and, in particular, Bradley-Terry) models admits a representation of the form ${r(x, y) = \beta \log \frac{\pi(y\mid x)}{\pi_\text{ref}(y\mid x)}}$ for an appropriate model $\pi(y\mid x)$ and a fixed reference policy $\pi_\text{ref}(y \mid x)$.
\end{theorem}

\begin{sproof}
    Take any reward function $r(x, y)$ and let its induced optimal model be $\pi_r(y \mid x)$ as given in Eq. \ref{eq:op_policy}. We demonstrate that within the equivalence class containing $r$, there exists a reward function with the specified reparameterized form. Define the following projection:
\begin{equation}
    f(r; \pi_\text{ref}, \beta)(x, y) = r(x, y) - \beta\log\sum_{y}\pi_\text{ref}(y\mid x)\exp\left(\frac{1}{\beta}r(x, y)\right)
\end{equation}
This operator $f$ rescales the reward by subtracting a normalization factor involving the logarithm of the partition function over $y$ for each $x$. Importantly, since the normalization term is solely a function of $x$, $f(r; \pi_\text{ref}, \beta)(x, y)$ remains in the equivalence class of $r(x, y)$. If we now substitute $r$ with the right-hand side of Eq.~\ref{eq:main_eq} (a relation that holds universally for reward functions), we see that $f(r; \pi_\text{ref}, \beta)(x, y) = \beta \log \frac{\pi_r(y\mid x)}{\pi_\text{ref}(y\mid x)}$. Thus, the map $f$ produces a representative from the equivalence class with the desired form, confirming that our proposed reparameterization entails no loss of expressiveness.
\end{sproof}

Theorem~\ref{thm:main} can also be interpreted as identifying the specific reward function chosen by the DPO reparameterization within each equivalence class, namely, the one that fulfills:

\begin{equation}\label{eq:lag_p}
     \sum_{y}\underbrace{\pi_\text{ref}(y\mid x)\exp\left(\frac{1}{\beta}r(x, y)\right)}_{=\pi(y\mid x)\text{, using Thm.~\ref{thm:main} reparam.}} = 1,
\end{equation}

This means that $\pi(y\mid x)$ is a proper probability distribution, as all probabilities are nonnegative and total to 1. Furthermore, referring to Eq.~\ref{eq:op_policy}, Eq.~\ref{eq:lag_p} represents the partition function corresponding to the optimal policy determined by the reward $r(x, y)$. A central contribution of the DPO framework is to enforce particular constraints on the otherwise underdetermined Plackett-Luce family (and specifically the Bradley-Terry model) of preference-based models. These constraints ensure that the spectrum of reward models remains intact, while at the same time making the optimal policy in Eq.~\ref{eq:op_policy} straightforward to compute analytically for any prompt $x$.

\subsection{Instability of Actor-Critic Algorithms} Our approach also facilitates the analysis of instability issues found in common actor-critic methods for RLHF, such as PPO. Consistent with the RLHF methodology, we examine the RL fine-tuning stage described in Section~\ref{section:prelims}. This analysis highlights parallels with the control as inference perspective on constrained RL \cite{levine2018reinforcement} set forth in Eq.~\ref{eq:RL}. Suppose we have a policy model $\pi_{\theta}(y\mid x)$; we can minimize $\mathbb{D}_{\text{KL}}[\pi_{\theta}(y|x) \mid \mid \pi^*(y\mid x)]$, with $\pi^*$ denoting the optimal policy from Eq.~\ref{eq:optimum_model} as determined by the reward $r_{\phi}(y, x)$. After algebraic manipulation, the resulting optimization target becomes:

\begin{equation}\label{eq:AC}
    \max_{\pi_{\theta}}\mathbb{E}_{\pi_{\theta}(y\mid x)}\bigg[\underbrace{r_{\phi}(x, y) -\beta\log\sum_{y}\pi_\text{ref}(y\mid x)\exp\left(\frac{1}{\beta}r_{\phi}(x, y)\right)}_{f(r_{\phi}, \pi_\text{ref}, \beta)} - \underbrace{\beta\log\frac{\pi_{\theta}(y\mid x)}{\pi_\text{ref}(y\mid x)}}_{\text{KL}}\bigg]
\end{equation}

The objective under consideration aligns with that targeted in earlier studies 
\citep{ziegler2020finetuning, stiennon2022learning, bai2022training, ouyang2022training}, where the DPO-equivalent reward function is employed for the reward family $r_{\phi}$. Here, the normalization component in $f(r_{\phi}, \pi_\text{ref}, \beta)$ can be viewed as the soft value function associated with the reference policy $\pi_\text{ref}$. Although this normalization term does not influence the optimal policy, omitting it can lead to high variance in the policy gradient, thereby destabilizing the learning process. One strategy to manage the normalization is to introduce a parameterized value function; however, this approach introduces additional optimization challenges. An alternative adopted in prior works involves normalizing rewards based on human completion baselines, which serve as single-sample Monte Carlo approximations of the normalization term. In comparison, the DPO reparameterization offers a reward function formulation that obviates the need for such baseline estimates.

\section{Experiments}
This section presents an empirical study of DPO for training policies directly from human preferences. To begin, in a controlled environment for text generation, we investigate the sample efficiency of DPO in balancing the maximization of the reward signal against minimizing KL-divergence from a reference policy, in comparison to standard preference-based learning approaches such as PPO. Subsequently, we test DPO on more complex RLHF tasks and on larger model scales, including applications in summarization and dialogue. Our results indicate that, even with minimal hyperparameter tuning, DPO typically matches or surpasses strong baselines, such as PPO-based RLHF and the best-out-of-$N$ trajectory selection using a learned reward model. Prior to detailing these empirical findings, we outline our experimental methodology; further implementation specifics are provided in Appendix~\ref{app:exp_details}.

\textbf{Tasks.} We conduct experiments on three distinct open-ended text generation scenarios. For each setting, the algorithms are trained to learn a policy using a preference dataset $\mathcal{D}=\bigl\{x^{(i)}, y_w^{(i)}, y_l^{(i)}\bigr\}_{i=1}^N$. In the case of \textbf{controlled sentiment generation}, $x$ denotes the opening segment of a movie review sourced from the IMDb dataset \cite{maas-EtAl:2011:ACL-HLT2011}, with the goal being to generate a continuation $y$ exhibiting positive sentiment. To facilitate a controlled assessment, we automatically generate preference pairs for outputs by employing a pre-trained sentiment classifier such that $p(\text{positive}\mid x,y_w)>p(\text{positive}\mid x,y_l)$. For the SFT baseline, we further train GPT-2-large on the IMDb training set reviews until performance stabilizes (implementation details can be found in App~\ref{app:sentiment_details}). Within the \textbf{summarization} task, $x$ is a Reddit forum entry and the task for the policy is to produce a summary $y$ capturing the essential content of the post. As in prior research, we utilize the Reddit TL;DR summarization corpus \citep{volske-etal-2017-tl} and pair it with human preference data obtained by \citeauthor{stiennon2022learning}. Our SFT baseline is trained on summaries authored by humans for forum posts,\footnote{\url{https://huggingface.co/CarperAI/openai_summarize_tldr_sft}} and we leverage the TRLX library \citep{leandro_von_werra_2023_7790115} for RLHF training. Notably, the human preference annotations were collected using generations from a separate, albeit similarly fine-tuned, SFT model, as described by \citeauthor{stiennon2022learning}. The third scenario, \textbf{single-turn dialogue}, presents $x$ as a user-provided query that spans domains ranging from science inquiries to personal advice. Here, the policy must yield an informative and engaging reply $y$. For this, we adopt the Anthropic Helpful and Harmless dialogue dataset \citep{bai2022training}, which includes 170,000 dialogues between people and an AI assistant, where each interaction concludes with two model-generated responses and a human-annotated preference label. Due to the absence of a ready-made SFT baseline in this scenario, we fine-tune a pre-existing language model exclusively on the responses marked as preferred, thereby constructing the SFT model.

\textbf{Evaluation.} We adopt two main strategies for evaluating our models. To systematically assess how well each algorithm optimizes for constrained reward maximization, we examine their performance in controlled sentiment generation. Here, we plot each algorithm’s trade-off between attained reward and KL-divergence relative to the reference policy, leveraging our access to the oracle reward function (provided by a sentiment classifier) to construct this frontier. Conversely, since direct reward ground truth is inaccessible in practical scenarios, we use the \textit{win rate} over a baseline policy, gauged with GPT-4 as an automated surrogate for human judgment of summary quality and dialogue helpfulness in the summarization and single-turn dialogue settings, respectively. For the summarization task, we set the test set reference summaries as our baseline, while in the dialogue setting, we compare against the test dataset's preferred responses. While literature indicates that language models often outperform conventional metrics as evaluators \citep{Chen2023ExploringTU}, we also run a human evaluation to support the use of GPT-4 as a judge, as detailed in Sec.~\ref{sec:human-judgments}. Our findings indicate a strong correlation between GPT-4’s ratings and those of human evaluators, with GPT-4’s alignment with human preferences meeting or exceeding typical inter-annotator agreement.

\textbf{Methods.} Besides DPO, our experimental benchmarks include several established methods for aligning language model outputs with human preferences. We apply zero-shot prompting with \textbf{GPT-J} \citep{gpt-j} for summarization, and adopt 2-shot prompting using \textbf{Pythia-2.8B} \citep{biderman2023pythia} for single-turn dialogue. We further examine the \textbf{SFT} approach, as well as \textbf{Preferred-FT}, a model supervised to fit the selected completion $y_w$ generated either from SFT (for sentiment and summarization) or from a general language model (for dialogue). An additional approach we consider is \textbf{Unlikelihood}~\citep{welleck2019neural}, which involves maximizing the likelihood of $y_w$ and explicitly \textit{minimizing} that of $y_l$, governed by a tunable `unlikelihood' weight $\alpha\in[0,1]$. Our comparisons also include \textbf{PPO} \citep{schulman2017proximal}, utilizing a reward model trained on preference annotations, as well as \textbf{PPO-GT}, an oracle variant directly optimizing with access to the ground-truth reward in the controlled sentiment context. In the sentiment experiments, we use two PPO-GT implementations: a standard one from \cite{leandro_von_werra_2023_7790115} and a variant with reward normalization and additional hyperparameter tuning for improved performance—these enhancements are also applied to PPO runs with learned rewards. Finally, we evaluate a \textbf{Best of $N$} baseline, which involves sampling $N$ outputs from the SFT model (or Preferred-FT for dialogue) and selecting the top candidate as determined by a reward model trained on preference data. While this approach can achieve high quality, it is computationally burdensome for even moderate $N$, due to the need to generate $N$ outputs per input at test time.

\subsection{How well can DPO optimize the RLHF objective?}

The RLHF framework typically employs a KL-constrained reward maximization objective, which serves to balance maximizing reward with preventing significant divergence from a reference policy. Consequently, evaluating different algorithms necessitates considering both their attained reward and the magnitude of KL divergence; an algorithm yielding higher reward at the expense of much greater KL may not be preferable. Figure~\ref{fig:frontier-tldr-main} presents the trade-off curves between reward and KL for a variety of methods in the sentiment analysis context. To examine these trade-offs, we conduct several independent training experiments per algorithm, varying the conservativeness hyperparameter in each instance (target KL $\in\{3,6,9,12\}$ for PPO; $\beta \in \{0.05,0.1,1,5\}$ for DPO; $\alpha\in\{0.05,0.1,0.5,1\}$ for unlikelihood; multiple random seeds for preferred-FT), resulting in a total of 22 training runs. Throughout training, at intervals of 100 steps up to convergence, policies are evaluated on a set of test prompts. For each evaluation, we measure the mean reward based on the true reward function and the mean sequence-level KL divergence\footnote{Specifically, we aggregate the KL-divergence computed at each timestep across a sequence.} relative to the reference policy, $\text{KL}\left(\pi\mid \mid \pi_\text{ref}\right)$. The results indicate that DPO delivers the most favorable reward-KL trade-off, attaining superior reward without incurring excessive KL. This outcome is significant for several reasons. Notably, although DPO and PPO both target the same objective, DPO exhibits far greater efficiency, offering a reward/KL profile that unequivocally surpasses that of PPO. Furthermore, DPO is able to surpass PPO even in scenarios where PPO is supplied with access to ground truth rewards (PPO-GT).

\subsection{Can DPO scale to real preference datasets?}
\label{sec:dpo-real-datasets}

We next investigate the fine-tuning effectiveness of DPO for tasks involving summarization and single-turn dialogue. In summarization, standard automated metrics such as ROUGE have been shown to have weak correlation with human judgment~\citep{stiennon2022learning}. Previous studies have demonstrated that preference-based fine-tuning with PPO enhances the quality of language model-generated summaries. To assess method performance, we generate completions for the test portion of the TL;DR summarization dataset and measure the mean win rate when compared to reference outputs from the test set. For all evaluated methods, completions are sampled across a range of temperatures between 0.0 and 1.0, with corresponding win rates illustrated in Figure~\ref{fig:frontier-tldr-main} (right). DPO, PPO, and Preferred-FT are all applied to fine-tune an identical GPT-J SFT model\footnote{\url{https://huggingface.co/CarperAI/openai_summarize_tldr_sft}}. Our results show that DPO achieves an approximate 61\% win rate at temperature 0.0, outperforming PPO, which reaches a win rate near 57\% at its optimal temperature setting. Additionally, DPO attains a superior maximum win rate relative to the best-of-$N$ baseline. It is important to mention that no significant hyperparameter optimization was conducted for DPO's $\beta$, suggesting that DPO's full capabilities may not be captured here. Furthermore, DPO displays greater stability with respect to sampling temperature in comparison to PPO, which can regress to baseline GPT-J performance at elevated temperatures. The Preferred-FT approach does not demonstrate notable improvement over the initial SFT model. We further perform direct human evaluations contrasting DPO and PPO in Section~\ref{sec:human-judgments}, observing that DPO samples drawn at a temperature of 0.25 are favored 58\% of the time over PPO samples at the same temperature.

For single-turn dialogue tasks, we assess various approaches using the subset of the Anthropic HH dataset’s test split \citep{bai2022training} that features a single human-assistant exchange. To measure win rates, GPT-4 judges use the test set’s preferred completions as references for each method. Since there is no canonical SFT baseline for this setting, we begin with a pre-trained Pythia-2.8B model. We employ Preferred-FT to fit a reference model on the preferred completions, ensuring the outputs remain in-distribution, and subsequently apply DPO training. Additionally, we benchmark against the strongest result from 128 Preferred-FT generations (noting that the Best of $N$ strategy reaches saturation at $N=128$ for this dataset; refer to Appendix Figure~\ref{fig:best-of-n}), as well as against a two-shot prompted Pythia-2.8B base model. DPO matches or outperforms these baselines when using their optimal sampling temperatures. We further include an RLHF system trained via PPO on the Anthropic HH dataset\footnote{\url{https://huggingface.co/reciprocate/ppo_hh_pythia-6B}} developed by a prominent group\footnote{\url{https://github.com/CarperAI/trlx/tree/main/examples/hh}}; however, across different prompting and sampling schemes, its performance does not exceed that of the unmodified Pythia-2.8B. Drawing from our TL;DR findings and the commonality of reward objectives, we treat the Best of 128 baseline as a loose stand-in for PPO-level outcomes. In sum, DPO is uniquely efficient and reliably outperforms the set of preferred completions in the Anthropic HH dataset, offering comparable or improved results relative to the much more computationally intensive Best of 128 approach. As illustrated in Figure~\ref{fig:dialogue-main}, DPO also attains its peak effectiveness in only a few training steps.

\subsection{Generalization to a new input distribution}

To deepen our understanding of how PPO and DPO respond to distributional shifts, we assess the policies learned in our Reddit TL;DR summarization setting on a distinct data distribution: news articles from the test set of the CNN/DailyMail corpus \citep{nallapati-etal-2016-abstractive}. In this evaluation, we use the optimal sampling temperatures from TL;DR (0 and 0.25). Table~\ref{tab:ood} shows the outcomes. For this comparison, we measured the win rate of GPT-4 over the human-written reference summaries for each dataset, adapting the same GPT-4 (C) evaluation prompt as for Reddit TL;DR, but substituting ``forum post'' with ``news article.'' Under this shift to a novel input distribution, DPO again substantially outperforms the PPO policy. These findings suggest that, despite lacking access to the unlabeled Reddit TL;DR prompts utilized by PPO, DPO-trained policies can achieve generalization on par with, or superior to, PPO-trained counterparts.

\subsection{Validating GPT-4 judgments with human judgments}
\label{sec:human-judgments}
We conduct a study involving human annotators to evaluate how well GPT-4's assessments correspond to human preferences, leveraging outcomes from the TL;DR summarization experiments and employing two types of GPT-4 prompts. The \textbf{GPT-4 (S)} (simple) prompt asks which summary most effectively conveys the essential information from the post. The \textbf{GPT-4 (C)} (concise) prompt adds a request to judge conciseness, since we observed that GPT-4, when prompted with the (S) version, often favors lengthier and more redundant summaries than those preferred by human evaluators. The full prompts are included in Appendix~\ref{app:prompts}. To cover a range of output qualities, we evaluate three methods: the best-performing DPO (temp. 0.25), the lowest-performing PPO (temp. 1.0), and an intermediate SFT method (temp. 0.25), each compared against PPO with greedy sampling (its optimal temperature). Our findings indicate that, across both prompts, GPT-4's agreement with human judgments approximates the level of consistency observed between human raters themselves; given that human rater availability was limited, we gathered repeated human annotations only for the DPO and PPO-1 scenarios. Overall, the \textbf{GPT-4 (C)} prompt yields win rates more closely aligned with human preferences, motivating its use in the primary analysis in Section~\ref{sec:dpo-real-datasets}. Further methodological specifics, such as the rater interface and participant roster, are detailed in Appendix~\ref{app:human-study}.

\section{Discussion}
Preference-based learning offers a robust and flexible approach for developing competent and aligned language models. In this work, we presented DPO, a streamlined method for leveraging preferences to train language models that bypasses the need for reinforcement learning. Instead of reformulating preference learning to fit within the conventional RL framework to utilize pre-existing RL methods, DPO constructs a correspondence between language model policies and reward structures. This enables the direct optimization of language models towards human preferences using a straightforward cross-entropy loss, eliminating both the necessity of reinforcement learning and any compromise in generality. DPO matches or exceeds the performance of established RLHF methods, including those built on PPO, even in the absence of significant hyperparameter tuning; as a result, DPO substantially lowers the practical challenges involved in preference-based language model training.

\textbf{Limitations \& Future Work.} The findings here point towards several avenues for subsequent research. A key question concerns the ability of the DPO policy to generalize to out-of-distribution scenarios, particularly in contrast to policies trained using explicit reward signals. Our preliminary experiments indicate that DPO's generalization is on par with that of PPO-based approaches, though a more thorough evaluation is warranted. It is also of interest to investigate whether self-labeling from a DPO policy can similarly exploit unannotated prompts effectively. Additionally, further investigation is needed into the dynamics of reward over-optimization in the direct preference optimization framework, including whether the modest drop in performance observed in Figure~\ref{fig:dialogue-main}-right constitutes an instance of this issue. While our experiments covered models with up to 6B parameters, scaling DPO to much larger, cutting-edge models remains an open and promising area. In terms of assessment, our analysis reveals that GPT-4-derived win rates can vary with the prompt, suggesting a need for research into the most reliable techniques for eliciting robust automated judgments. Lastly, the DPO approach shows promise in a range of generative modeling tasks beyond language model alignment through human preferences, pointing to numerous potential future applications.